\BourbakiAppendixExercises{习题}

\begin{enumerate}
\def\labelenumi{\arabic{enumi})}
\tightlist
\item
  设 K 是一个交换域，E 是 K 上的一个向量空间。我们用 \(\operatorname{End}_{\mathrm{K}}^{\mathrm{pf}}(\mathrm{E})\) 表示 \(\operatorname{End}_{\mathrm{K}}(\mathrm{E})\) 中由某个幂属于 \(\operatorname{End}_{\mathrm{K}}^{\mathrm{f}}(\mathrm{E})\) 的自同态组成的子集。若 \(u \in \operatorname{End}_{\mathrm{K}}^{\mathrm{pf}}(\mathrm{E})\)，则置 \(I_{u} = \bigcap_{n \geqslant 0} Im u^{n}\) 与 \(\operatorname{Tr} u = \operatorname{Tr}(u | I_{u})\)；对有限秩自同态，这个概念与第 2 目中定义的一致。a) 设 \(u \in \operatorname{End}_{\mathrm{K}}^{\mathrm{pf}}(\mathrm{E})\)；设 V 是 E 的一个在 u 下稳定的子空间，\(u_{V}\) 与 \(u_{E/V}\) 是由 u 诱导的 V 与 E/V 的自同态。证明等式 \(\operatorname{Tr} u = \operatorname{Tr} u_{V} + \operatorname{Tr} u_{E/V}\)。
\end{enumerate}

\begin{enumerate}
\def\labelenumi{\alph{enumi})}
\setcounter{enumi}{1}
\item
  设 U 是 \(\operatorname{End}_{\mathbf{K}}(\operatorname{E})\) 的一个子空间。假设存在一个整数 n，使得 U 中 n 个元素的每个乘积都有有限秩。证明映射 \(Tr: U \to K\) 是线性的。
\item
  设 F 是 K 上的一个向量空间，\(u: E \to F\) 与 \(v: F \to E\) 是使得 \(u \circ v \in \operatorname{End}_{\mathrm{K}}^{\mathrm{pf}}(\mathrm{F})\) 的同态。我们有 \(v \circ u \in \operatorname{End}_{\mathrm{K}}^{\mathrm{pf}}(\mathrm{E})\) 与 \(\operatorname{Tr} v \circ u = \operatorname{Tr} u \circ v\)。
\end{enumerate}

¶ 2) 保留上一个习题的记号。若 V 与 W 是 E 的子空间，则我们用 V \(\prec\) W 表示关系“W 在 V + W 中具有有限余维数”。

\begin{enumerate}
\def\labelenumi{\alph{enumi})}
\item
  设 V 是 E 的一个子空间。我们用 \(\mathfrak{A}\)（相应地 \(\mathfrak{a},\mathfrak{b}\)）表示 V 的自同态 \(u\) 中满足 \(u(\mathrm{V})\prec \mathrm{V}\)（相应地 \(u(\mathrm{E})\prec \mathrm{V}\)、\(u(\mathrm{V})\prec 0\)）者所成的集合。证明 \(\mathfrak{A}\) 是 \(\operatorname{End}_{\mathrm{K}}(\mathrm{E})\) 的一个 K-子代数，\(\mathfrak{a}\) 与 \(\mathfrak{b}\) 是 \(\mathfrak{A}\) 的两个双边理想，其和为 \(\mathfrak{A}\)，且我们有 \(\mathfrak{a} \cap \mathfrak{b} \subset \operatorname{End}_{\mathrm{K}}^{\mathrm{pf}}(\mathrm{E})\)。
\item
  设 \(u, u'\) 是 \(\mathfrak{A}\) 中可交换的元素；我们可以写 \(u = a + b\) 与 \(u' = a' + b'\)，其中 \(a, a'\) 属于 \(\mathfrak{a}\)，\(b, b'\) 属于 \(\mathfrak{b}\)。证明元素 \([a, a'] = aa' - a'a\)（\(\mathfrak{A}\) 中的）属于 \(\mathfrak{a} \cap \mathfrak{b}\)，且它的迹不依赖于 \(u\) 与 \(u'\) 的分解的选取；我们把它记作 \(\langle u, u' \rangle\)。
\item
  设 A 是 \(\mathfrak{A}\) 的一个交换子代数。证明存在唯一的 K-线性映射 \(\operatorname{Res}:\Omega_{\mathrm{K}}(\mathrm{A})\to \mathrm{K}\)，使得 \(\operatorname{Res}(udv) = \langle u,v\rangle\)（对 A 中所有的 \(u,v\)）（利用恒等式 \([u,vw] + [v,wu] + [w,uv] = 0\)）。
\item
  设 \(u, v \in A\)；假设 \(v(V) \subset V\)。设 \(\pi\) 是 E 的以 V 为像的投影。证明等式 \(\operatorname{Res}(u dv) = \operatorname{Tr}[\pi u, v]\)。由此推出，若 \(u(V) \subset V\)，则我们有 \(\operatorname{Res} u dv = 0\)（注意自同态 \([\pi u, v]\) 的平方为零）。
\item
  设 \(u \in \mathrm{A}\)，\(n \in \mathbf{N}\)。证明我们有 \(\operatorname{Res} u^n du = 0\)。若 \(u\) 在 \(\mathrm{A}\) 中可逆，则我们有 \(\operatorname{Res} u^{-n} du = 0\)（对 \(n \geqslant 2\)）；若此外 \(u(\mathrm{V}) \subset \mathrm{V}\)，则我们有 \(\operatorname{Res} u^{-1} du = \dim_{\mathrm{K}} \mathrm{V} / u(\mathrm{V})\)。
\end{enumerate}

\(f)\) 取 \(\mathrm{E} = \mathrm{K}((\mathrm{X}))\)（IV, §4, No.~9, 第 38 页），\(\mathrm{V} = \mathrm{K}[[\mathrm{X}]]\)，并让 \(\mathrm{A} = \mathrm{E}\) 通过位似作用在自身上。证明 \(c)\) 的假设被满足。设 \(f(\mathrm{X}) = \sum_{n} a_{n} \mathrm{X}^{n} d\mathrm{X}\) 是 \(\mathrm{K}((\mathrm{X}))\) 的一个元素。证明 \(\operatorname{Res} f(\mathrm{X}) d\mathrm{X}\) 等于 \(a_{-1}\)。由此推出：若 \(g\) 是 \(\mathrm{K}[[\mathrm{X}]]\) 的一个元素，满足 \(g(0) = 0\) 且 \(g'(0) \neq 0\)，则 \(\mathrm{X}^{-1}\) 在 \(f(\mathrm{X})\) 与 \(f(g(\mathrm{X})) g'( \mathrm{X})\) 的形式级数展开中的系数相同。

\backmatter
